\begin{longtable}[l]{ |m{8.2cm}|r|c|r|m{9.3cm}| }    
\caption{BLOCK0 参数}
    \label{tab:efuse-block0-para} \\\hline
    
    \rowcolor{lightgray}
        \textbf{参数} & \textbf{\thead{位宽}} & \textbf{\thead{硬件使用}} & \textbf{\thead{\hyperref[fielddesc:EFUSEWRDIS]{EFUSE\_WR\_DIS} \\烧写保护位}}
    & \textbf{描述} \\
    \endfirsthead
    
    \multicolumn{5}{c}{\bfseries \tablename\ \thetable{} -- 接上页}\\\hline
    
    \rowcolor{lightgray}
        \textbf{参数} & \textbf{\thead{位宽}} & \textbf{\thead{硬件使用}} & \textbf{\thead{\hyperref[fielddesc:EFUSEWRDIS]{EFUSE\_WR\_DIS} \\烧写保护位}}
    & \textbf{描述} \\
    \endhead 
    \multicolumn{5}{r}{\textbf{见下页}} \\
    \endfoot
    \endlastfoot
    \hline   
    \hyperref[fielddesc:EFUSEWRDIS]{\textbf{EFUSE\_WR\_DIS}} & 32 & Y & N/A & 表示是否禁止 eFuse 烧写 \\
    \hline
    \hyperref[fielddesc:EFUSERDDIS]{\textbf{EFUSE\_RD\_DIS}}  & 7 & Y & 0  & 表示是否禁止用户读取 eFuse BLOCK4 \verb+~+ 10 的内容 \\
    \hline
    \hyperref[fielddesc:EFUSEDISICACHE]{EFUSE\_DIS\_ICACHE}  &  1 & Y & 2  & 表示是否禁用 ICache \\
    \hline
    \hyperref[fielddesc:EFUSEDISDCACHE]{EFUSE\_DIS\_DCACHE}  & 1 & Y & 2  & 表示是否禁用 DCache \\
    \hline
    \hyperref[fielddesc:EFUSEDISDOWNLOADICACHE]{EFUSE\_DIS\_DOWNLOAD\_ICACHE}  & 1 & Y & 2 & 表示是否在下载模式下关闭 ICache\\   
    \hline
    \hyperref[fielddesc:EFUSEDISDOWNLOADDCACHE]{EFUSE\_DIS\_DOWNLOAD\_DCACHE} & 1 & Y & 2 & 表示是否在下载模式下关闭 DCache\\      
    \hline
    \hyperref[fielddesc:EFUSEDISFORCEDOWNLOAD]{EFUSE\_DIS\_FORCE\_DOWNLOAD} & 1 & Y & 2 & 表示是否禁止强制芯片进入下载模式\\   
    \hline
    \hyperref[fielddesc:EFUSEDISUSBOTG]{EFUSE\_DIS\_USB\_OTG} &1 & Y & 2 & 表示是否禁用 USB OTG 功能 \\     
    \hline
    \hyperref[fielddesc:EFUSEDISTWAI]{EFUSE\_DIS\_TWAI} & 1 & Y & 2 & 表示是否禁用 TWAI 控制器功能 \\  
    \hline
    \hyperref[fielddesc:EFUSEDISAPPCPU]{EFUSE\_DIS\_APP\_CPU} & 1 & Y & 2 & 表示是否禁用 app CPU \\   
    \hline
    \hyperref[fielddesc:EFUSESOFTDISJTAG]{EFUSE\_SOFT\_DIS\_JTAG} & 3 & Y & 31 & 表示是否软关断 JTAG 功能\\    
    \hline
    \hyperref[fielddesc:EFUSEDISPADJTAG]{EFUSE\_DIS\_PAD\_JTAG} & 1 & Y & 2 & 表示是否永久禁用 JTAG 功能\\   
    \hline
    \hyperref[fielddesc:EFUSEDISDOWNLOADMANUALENCRYPT]{EFUSE\_DIS\_DOWNLOAD\_MANUAL\_ENCRYPT} &  1 & Y & 2 & 表示是否在 download boot 模式下禁用 flash 加密功能\\ 
    \hline
    \hyperref[fielddesc:EFUSEUSBEXCHGPINS]{EFUSE\_USB\_EXCHG\_PINS} & 1 & Y & 30 & 表示是否交换 USB D+/D- 管脚 \\
    \hline
    \hyperref[fielddesc:EFUSEEXTPHYENABLE]{EFUSE\_EXT\_PHY\_ENABLE} & 1 & N & 30 & 表示是否使能外部 USB PHY \\
    \hline
    \hyperref[fielddesc:EFUSEVDDSPIXPD]{EFUSE\_VDD\_SPI\_XPD} & 1 & Y & 3 & 表示 flash 电压调节器是否上电 \\        
    \hline
    \hyperref[fielddesc:EFUSEVDDSPITIEH]{EFUSE\_VDD\_SPI\_TIEH} & 1 & Y & 3 & 表示 flash 电压调节器是否短接至 VDD\_RTC\_IO \\   
    \hline
    \hyperref[fielddesc:EFUSEVDDSPIFORCE]{EFUSE\_VDD\_SPI\_FORCE} & 1 & Y & 3 & 表示是否强制使用 EFUSE\_VDD\_SPI\_XPD 和 EFUSE\_VDD\_SPI\_TIEH 配置 flash 电压 LDO \\ 
    \hline
    \hyperref[fielddesc:EFUSEWDTDELAYSEL]{EFUSE\_WDT\_DELAY\_SEL} & 2 & Y & 3 & 表示 RTC 看门狗超时阈值 \\
    \hline
    \hyperref[fielddesc:EFUSESPIBOOTCRYPTCNT]{EFUSE\_SPI\_BOOT\_CRYPT\_CNT} & 3 & Y & 4 & 表示是否使能 SPI boot 加解密 \\
    \hline
    \hyperref[fielddesc:EFUSESECUREBOOTKEYREVOKE0]{EFUSE\_SECURE\_BOOT\_KEY\_REVOKE0} & 1 & N & 5 & 表示是否使能撤销第一个安全启动密钥 \\
    \hline
    \hyperref[fielddesc:EFUSESECUREBOOTKEYREVOKE1]{EFUSE\_SECURE\_BOOT\_KEY\_REVOKE1} & 1 & N & 6 & 表示是否使能撤销第二个安全启动密钥 \\
    \hline
    \hyperref[fielddesc:EFUSESECUREBOOTKEYREVOKE2]{EFUSE\_SECURE\_BOOT\_KEY\_REVOKE2} & 1 & N & 7 & 表示是否使能撤销第三个安全启动密钥 \\
    \hline
    \hyperref[fielddesc:EFUSEKEYPURPOSE0]{EFUSE\_KEY\_PURPOSE\_0} & 4 & Y & 8 & 表示 Key0 用途，见表 \ref{tab:efuse-key-purpose} \\
    \hline
    \hyperref[fielddesc:EFUSEKEYPURPOSE1]{EFUSE\_KEY\_PURPOSE\_1} & 4 & Y & 9 & 表示 Key1 用途，见表 \ref{tab:efuse-key-purpose} \\
    \hline
    \hyperref[fielddesc:EFUSEKEYPURPOSE2]{EFUSE\_KEY\_PURPOSE\_2} & 4 & Y & 10 & 表示 Key2 用途，见表 \ref{tab:efuse-key-purpose} \\ 
    \hline
    \hyperref[fielddesc:EFUSEKEYPURPOSE3]{EFUSE\_KEY\_PURPOSE\_3} & 4 & Y & 11 & 表示 Key3 用途，见表 \ref{tab:efuse-key-purpose} \\    
    \hline
    \hyperref[fielddesc:EFUSEKEYPURPOSE4]{EFUSE\_KEY\_PURPOSE\_4} & 4 & Y & 12 & 表示 Key4 用途，见表 \ref{tab:efuse-key-purpose} \\    
    \hline
    \hyperref[fielddesc:EFUSEKEYPURPOSE5]{EFUSE\_KEY\_PURPOSE\_5} & 4 & Y & 13 & 表示 Key5 用途，见表 \ref{tab:efuse-key-purpose} \\  
    \hline
    \hyperref[fielddesc:EFUSESECUREBOOTEN]{EFUSE\_SECURE\_BOOT\_EN} & 1 & N & 15 & 表示是否使能安全启动 \\ 
    \hline
    \hyperref[fielddesc:EFUSESECUREBOOTAGGRESSIVEREVOKE]{EFUSE\_SECURE\_BOOT\_AGGRESSIVE\_REVOKE} & 1 & N & 16 & 表示是否使 Secure Boot 的撤销采用激进策略 \\
    \hline
    \hyperref[fielddesc:EFUSEDISUSBJTAG]{EFUSE\_DIS\_USB\_JTAG} & 1 & Y & 2 & 表示是否禁用 usb\_serial\_jtag 模块的 usb 转 jtag 功能 \\
    \hline
    \hyperref[fielddesc:EFUSEDISUSBSERIALJTAG]{EFUSE\_DIS\_USB\_SERIAL\_JTAG} & 1 & Y & 2 & 表示是否禁用 usb\_serial\_jtag 模块 \\
    \hline
    \hyperref[fielddesc:EFUSESTRAPJTAGSEL]{EFUSE\_STRAP\_JTAG\_SEL} & 1 & Y & 2 & 表示是否使能使用 strapping GPIO3 选择 usb\_to\_jtag 或 pad\_to\_jtag 的功能 \\
    \hline
    \hyperref[fielddesc:EFUSEUSBPHYSEL]{EFUSE\_USB\_PHY\_SEL} & 1 & Y & 2 & 表示内部 PHY、外部 PHY 和 USB OTG、USB Serial/JTAG 之间的连接关系 \\
    \hline
    \hyperref[fielddesc:EFUSEFLASHTPUW]{EFUSE\_FLASH\_TPUW} & 4 & N & 18 & 表示配置上电后 flash 等待时间 \\    
    \hline    
    \hyperref[fielddesc:EFUSEDISDOWNLOADMODE]{EFUSE\_DIS\_DOWNLOAD\_MODE} & 1 & N & 18 & 表示是否关闭所有 Download boot 模式\\       
    \hline
    \hyperref[fielddesc:EFUSEDISLEGACYSPIBOOT]{EFUSE\_DIS\_LEGACY\_SPI\_BOOT} & 1 & N & 18 & 表示是否关闭 Legacy SPI boot 模式 \\
    \hline
    \hyperref[fielddesc:EFUSEDISUSBPRINT]{EFUSE\_DIS\_USB\_PRINT} & 1 & N & 18 & 表示是否关闭 USB 打印 \\
    \hline 
    \hyperref[fielddesc:EFUSEFLASHECCMODE]{EFUSE\_FLASH\_ECC\_MODE} & 1 & N & 18 & 表示 ROM 中的 flash ECC 模式 \\
    \hline 
    \hyperref[fielddesc:EFUSEDISUSBSERIALJTAGDOWNLOADMODE]{EFUSE\_DIS\_USB\_SERIAL\_JTAG\_DOWNLOAD\_MODE} & 1 & N & 18 & 表示是否禁用 USB-Serial-JTAG 下载功能 \\
    \hline 
    \hyperref[fielddesc:EFUSEENABLESECURITYDOWNLOAD]{EFUSE\_ENABLE\_SECURITY\_DOWNLOAD} & 1 & N & 18 & 表示是否使能 UART 安全下载模式 \\
    \hline
    \hyperref[fielddesc:EFUSEUARTPRINTCONTROL]{EFUSE\_UART\_PRINT\_CONTROL} & 2 & N & 18 & 表示 UART boot 信息的默认打印方式 \\
    \hline
    \hyperref[fielddesc:EFUSEPINPOWERSELECTION]{EFUSE\_PIN\_POWER\_SELECTION} & 1 & N & 18 & 表示 GPIO33 \verb+~+ GPIO37、GPIO47、GPIO48 的供电来源 \\
    \hline
    \hyperref[fielddesc:EFUSEFLASHTYPE]{EFUSE\_FLASH\_TYPE} & 1 & N & 18 & 表示 SPI flash 的最大行数 \\
    \hline
    \hyperref[fielddesc:EFUSEFLASHPAGESIZE]{EFUSE\_FLASH\_PAGE\_SIZE} & 2 & N & 18 & 表示 flash 的页大小 \\
    \hline
    \hyperref[fielddesc:EFUSEFLASHECCEN]{EFUSE\_FLASH\_ECC\_EN} & 1 & N & 18 & 表示是否在 flash boot 模式下使能 ECC 功能 \\
    \hline
    \hyperref[fielddesc:EFUSEFORCESENDRESUME]{EFUSE\_FORCE\_SEND\_RESUME}  & 1 & N & 18 & 表示是否强制 ROM 代码在 SPI 启动过程中向 SPI flash 发送恢复命令 \\  
    \hline
    \hyperref[fielddesc:EFUSESECUREVERSION]{EFUSE\_SECURE\_VERSION}& 16 & N & 18 & 表示 IDF 安全版本 \\\hline
    \hyperref[fielddesc:EFUSEDISUSBOTGDOWNLOADMODE]{EFUSE\_DIS\_USB\_OTG\_DOWNLOAD\_MODE}  & 1 & N & 19 & 表示是否禁用 USB-OTG 下载功能 \\\hline
\end{longtable}
